\documentclass[11pt,a4paper]{article}
\usepackage[margin=25mm]{geometry}
\usepackage{amsmath,amssymb,amsthm,hyperref}
\setlength{\emergencystretch}{2em}
\newtheorem{lemma}{Lemma}
\newtheorem{theorem}[lemma]{Theorem}
\newcommand{\gen}[1]{\langle#1\rangle}
\title{An overlap bound between two unique crossings}
\author{Human Compression project: L044 continuation}
\date{}
\begin{document}\maketitle
Throughout $G=\gen{\rho_0,\ldots,\rho_{r-1}}\leq A_n$,
$n\geq12$, is transitive with a connected string C-marking,
and every maximal omission $G_k$ is intransitive. All actions
are the given faithful permutation actions. An eligible edge of
colour $k$ crosses distinct $G_k$-orbits.

\begin{lemma}[The outside orbit alone]\label{outside}
Suppose the unique eligible $i$-edge joins $a\in A$ to $b\in B$,
the two $G_i$-orbits, with $|B|\geq2$, and $G_{<i}$ fixes $B$ pointwise.
If $h>i$ and $\rho_h$ acts nontrivially on $A$, then
\[
                         |A|\geq i+h+1.                 \tag{1}
\]
For $i>0$ the inequality is strict unless $h=i+1$.
The conclusion uses no transitivity assumption on a final tail.
\end{lemma}
\begin{proof}
For $i=0$, the accepted fixed-prefix/no-gap lemma supplies a
simple path in $A$ from $a$ including a final $h$-edge, with
at least $h$ edges. Hence $|A|\geq h+1$.

Let $i>0$ and put $L=G_{<i}$. If $L$ is transitive on $A$,
every generator of index at least $i+2$ fixes $a$ and centralizes
$L$, so fixes $A$ pointwise. Thus $h=i+1$. The accepted
transitive-side bound gives $|A|\geq2(i+1)$, exactly (1).

Suppose $L$ is intransitive on $A$. Its commuting-orbit
normalization gives faithful equal-sized orbits of size
$k\geq i+1$; $\rho_{i+1}$ exchanges the anchor orbit
$\beta=a^L$ with another orbit $\beta'$. Put $d=h-i$.
For $d=1$, the two orbits give $|A|\geq2k\geq2i+2$.
For $d=2$, the first active orbit for $\rho_h$ is $\beta'$.
If it is moved to another orbit, there are at least three
orbits and $|A|\geq3k\geq3i+3>2i+3$.
If it is preserved nontrivially, the accepted internal-block
lemma gives $A=\beta\cup\beta'$ and $i\leq k/2-1$.
Consequently $|A|=2k\geq4i+4>2i+3$.
The internal-block lemma's nontrivial-$B$ hypothesis is explicit here.

For $d\geq3$, the accepted active-orbit distance lemma gives
at least $d$ orbits. Therefore
\[
 |A|\geq d(i+1)=2i+d+i(d-2)\geq2i+d+1=i+h+1.
\]
Equality is possible only for $i=1,d=3,k=2$ with exactly three
$L$-orbits. But then $L=\gen{\rho_0}$, and $\rho_0$ is a
transposition on each of those three pairs and fixes $B$.
It is odd, contrary to $G\leq A_n$. Thus this case is strict.
\end{proof}

\begin{lemma}[Three regions and the middle path]\label{middle}
Let $i<j$ be colours with unique eligible edges. Suppose the
$j$-edge lies in the $G_i$-orbit $B$ and the $i$-edge lies in
the other side of the $j$-cut. Write the two edges as
$a\mathrel{-_i}b$ and $c\mathrel{-_j}d$. Deleting them from
the full permutation graph gives three regions $A,M,D$,
with $a\in A$, $b,c\in M$, $d\in D$. If $j\geq i+2$, then
\[
                              |M|\geq j-i.              \tag{2}
\]
\end{lemma}
\begin{proof}
A unique eligible edge is a bridge of the full permutation
graph: all other edges preserve the two omission orbits,
and that edge is their only connection. The full graph is
connected by transitivity. Removing the two distinct bridges
therefore gives exactly three connected components in the
stated order. Here $A$ is the first $G_i$-orbit and $D$ is
the $G_j$-orbit not containing the $i$-edge.

The endpoints $b,c$ are distinct, since $\rho_j$ fixes $b$
by the unique $i$-crossing lemma but moves $c$.
Take a simple path in $M$ from $b$ to $c$. Suppose it misses
some colour $k$ with $i<k<j$. Its word factors as $uv$,
where $u$ uses indices less than $k$ and $v$ uses indices
greater than $k$; these factors commute by the string relations.
Every generator in $v$ fixes $b$, since its index is at least
$i+2$. Every generator in $u$ fixes $c$, since its index is
at most $j-2$. Thus, using right actions,
$c=b^{uv}=b^{vu}=b^u$, while $c^u=c$; applying $u^{-1}$
gives $b=c$, a contradiction. Each intermediate colour must
therefore occur. The path has at least $j-i-1$ edges and
$j-i$ vertices, proving (2).
\end{proof}

\begin{theorem}[The shared bands cannot strictly overlap at high rank]
\label{overlap}
Use the three regions of Lemma~\ref{middle}, with $j\geq i+2$.
Suppose $A,D$ are nontrivial, $G_{<i}$ fixes $M\cup D$,
and $G_{>j}$ fixes $A\cup M$. Let
\[
 i<h<j,\qquad 0\leq g<j,
\]
where $\rho_h$ acts nontrivially on $A$ and $\rho_g$ acts
nontrivially on $D$. Then
\[
                              n\geq2r+h-g.              \tag{3}
\]
In particular, if $r>\lfloor(n-1)/2\rfloor$, then $g\geq h$,
so also $g>i$.

If $g=h$ and the rank bound fails, then $n=2r$ and
\[
 |A|=i+h+1,\quad |M|=j-i,\quad |D|=2r-j-h-1.           \tag{4}
\]
Every simple middle path of minimum length is then Hamiltonian
on $M$ and uses each intermediate colour exactly once.
If $i>0$, equality forces $h=i+1$; if $j<r-1$, it forces
$h=j-1$. Thus when both crossings are interior the only surviving
index configuration is $j=i+2$, $h=g=i+1$, with
$|A|=2i+2$, $|M|=2$ and $|D|=2(r-1-j)+2$.
\end{theorem}
\begin{proof}
Apply Lemma~\ref{outside} to the first crossing. Apply its
literal reversal to the second crossing: the reversed indices
are $i'=r-1-j$ and $h'=r-1-g$, so
\[
 |A|\geq i+h+1,\qquad
 |D|\geq i'+h'+1=2r-j-g-1.
\]
The other omission orbit is nontrivial in both uses because
it contains $M$ and the opposite nontrivial outside region.
Adding (2) gives (3). If $g<h$, it gives $n\geq2r+1$,
contradicting failure of the rank target. If $g=h$, failure
forces equality everywhere and $n=2r$. This proves (4), the
path assertion, and both strictness consequences from
Lemma~\ref{outside}.
\end{proof}

\paragraph{Source routing and attribution.}
This is a new three-region count for the overlap issue in
Cameron--Fernandes--Leemans--Mixer, \emph{Highest rank of a polytope
for $A_n$}, arXiv:1605.09173v1, Proposition~5.28, physical page~37.
The accepted unique-crossing, no-gap, faithful commuting-orbit,
active-distance and internal-block arguments, including their
retained Whiston input, remain paid dependencies. The source's
claim $g>i$ follows here only after the stated orientations and
high-rank hypothesis, not just from membership in an active-label
set. No minimality or maximality of $g,h$ is needed for the count.

Equality $g=h$ is \emph{not} excluded. Neither are the separate
$g>h$ two-orbit equality case, the unproved source support routing,
or the complete later graph tables asserted recovered. No new
finite computation, formal declaration, original or question
credit is claimed.
\end{document}
